\chapter{Déploiement et maintenance}

\section{Environnements}
\begin{longtable}{p{4cm}p{10.5cm}}
\toprule \textbf{Environnement} & \textbf{URL} \\ \midrule \endhead
Aperçu (preview) & \url{https://id-preview--d3d48394-d747-4de9-844c-541a6a3403d6.lovable.app} \\
Production & \url{https://academieplus.lovable.app} \\
\bottomrule
\end{longtable}

Le projet Supabase associé porte la référence \texttt{lfothlxoixayjiytwwqa}.

\section{Variables d'environnement}
Le fichier \filepath{.env} est rempli automatiquement avec~:
\begin{lstlisting}[language=bash]
VITE_SUPABASE_URL=...
VITE_SUPABASE_PUBLISHABLE_KEY=...   # cle publique (anon) - OK cote client
VITE_SUPABASE_PROJECT_ID=...
\end{lstlisting}
La clé publique (\texttt{anon}) est volontairement publique~; la protection
réelle des données provient de la RLS. Les clés \emph{privées} restent dans les
secrets Supabase.

\section{Cycle de déploiement}
\begin{itemize}
  \item \textbf{Frontend} : build Vite (\texttt{npm run build}) servi en statique.
  \item \textbf{Edge Functions} : déploiement \textbf{automatique} à chaque
        modification (aucune commande manuelle).
  \item \textbf{Migrations} : les fichiers SQL de \filepath{supabase/migrations/}
        sont versionnés et appliqués via l'outil de migration.
\end{itemize}

\section{Internationalisation}
L'application est bilingue \textbf{français / arabe} via \texttt{i18next}
(\filepath{src/i18n/}). Les libellés sont stockés dans
\filepath{src/i18n/locales/fr.json} et \filepath{src/i18n/locales/ar.json}.
L'interface élève est majoritairement en \textbf{arabe} (RTL), tandis que les
espaces d'administration et éditorial sont en français.

\section{Tests}
\begin{itemize}
  \item \textbf{Unitaires / composants} : Vitest + Testing Library
        (\filepath{src/test/}). Le moteur de niveau (\filepath{src/lib/levelEngine.ts})
        étant constitué de fonctions pures, il est idéal à couvrir par des tests.
  \item \textbf{Bout en bout} : Playwright (\filepath{playwright.config.ts}).
\end{itemize}

\section{Conseils pour les prochaines équipes}
\begin{enumerate}
  \item \textbf{Logique de niveau} : toute évolution du calcul du niveau doit
        passer par \filepath{src/lib/levelEngine.ts} (fonctions pures, testables)
        plutôt que d'être dispersée dans les composants.
  \item \textbf{Accès données} : privilégier les services (\filepath{src/services})
        et hooks (\filepath{src/hooks}) existants plutôt que des appels Supabase
        directs dans les composants.
  \item \textbf{Sécurité} : ne jamais déplacer les rôles dans \texttt{profiles}~;
        toujours créer les \texttt{GRANT} et politiques RLS dans la même migration
        qu'un \texttt{CREATE TABLE}.
  \item \textbf{IA} : ajouter toute nouvelle intégration IA dans une Edge Function,
        jamais côté client.
  \item \textbf{Nettoyage} : la racine du dépôt contient des scripts ponctuels
        historiques (\filepath{fix*.mjs}, \filepath{patch*.cjs}\dots) qui ne font
        pas partie du runtime et peuvent être archivés.
\end{enumerate}

\section*{Conclusion}
\addcontentsline{toc}{section}{Conclusion}
AcadémiePlus est une plateforme adaptative complète couvrant l'ensemble du
parcours d'un élève~: positionnement initial, cours, exercices/quiz calibrés,
assistant IA, suivi de progression et reporting parental. Son architecture
(React + Supabase + Edge Functions IA) sépare clairement présentation, logique
métier pure et accès données, ce qui en facilite l'évolution.
